\scanpage{5}
% 1–6
Nous supposons maintenant $k$ de car. $0$, de sorte qu'on a la notion d'exponentielle d'un champ de vecteurs formel $\xi$ sur \deleted{\unsure{X_n|X}} un schéma formel lisse $X$,
\[
 (t,x)\mapsto(\exp t\xi).x\ :\ \underbrace{X_1}_{\text{``axe temps}}\times X\longrightarrow X
\]

% 7, left
% slanted
\margin{Intégration d'équations différentielles}

% 7–10
g) Si $X$ est un $A$-schéma formel (lisse) et $X$ un \inserted{$A$-}champ de vecteurs, alors la résolvante $X_1\times X\longrightarrow X$ est un $A$-morphisme.

% 11–12
(Bien entendu, il suffirait de le \deleted{faire}\inserted{demander} pour $X$ de la forme $X_n$\ldots).

% 13–16
Ceci implique \inserted{\unsure{et|est}} (par l'intro$\unsure{^{\mathrm n}}$ de la variable temps, devenant variable spatiale), \unsure{donc|son} équivaut, à

g') \unsure{Idem que|Comme|Analogue à} g), mais avec \deleted{``}un ``champ de vecteurs variable avec le temps''

% 17–20
$\xi(t)$\quad (
\diagram{\begin{tikzcd}[ampersand replacement=\&,column sep=small,row sep=small]
X_1\times X \arrow[rr] \arrow[dr,dashed,"\mathrm{pr}_2"'] \& \& \underbrace{\unsure{\mathbb T|T|\mathbf T}(X)}_{\substack{\text{variété }\deleted{\text{\unsure{actuelle|vectorielle|formelle}}}\text{\inserted{formelle}}\\\text{tangente de }X}} \arrow[dl,dashed] \\
\& X \&
\end{tikzcd}}
), en résolvant l'équation différentielle formelle
\[
 \frac{d}{dt}x(t)=\xi(t,x)
\]

% 21–23
La résolvante $x=\varphi(t,x_0)$ correspond à un $\varphi:X_1\times X\longrightarrow X$ qui est un $A$-morphisme.
% archive sheet number: 4
